\documentclass{article}
\usepackage[T5]{fontenc}
\usepackage[utf8]{inputenc}
\usepackage[vietnamese]{babel}
\usepackage[a4paper,margin=1in]{geometry}
\usepackage{tikz}
\usepackage{longtable}
\usepackage{float}
\usetikzlibrary{positioning,arrows.meta,shapes.geometric}

\title{Hybrid Encryption File Transfer System}
\author{Báo cáo hệ thống}
\date{\today}

\begin{document}
\maketitle

\section{Tổng quan hệ thống}
Hệ thống \textit{Hybrid Encryption File Transfer System} được xây dựng để truyền tệp an toàn giữa người gửi và người nhận. Mục tiêu chính là đảm bảo ba tính chất: bí mật dữ liệu, toàn vẹn nội dung, và khả năng xác thực tệp đã giải mã.

Mô hình mã hóa lai kết hợp hai kỹ thuật:
\begin{itemize}
	\item AES-256-GCM: mã hóa dữ liệu tệp với tốc độ cao, tạo thêm \textit{authentication tag} để kiểm tra toàn vẹn.
	\item RSA-OAEP (2048 bit): mã hóa khóa phiên AES để chỉ người nhận sở hữu private key mới giải được.
\end{itemize}
Cơ chế này giảm chi phí tính toán (so với mã hóa RSA trực tiếp cho tệp lớn) đồng thời giữ mức an toàn cao.

\section{Phân rã chức năng}
\subsection{Chức năng 1: Quản lý khóa (Key Management)}
\textbf{Mô tả:} Tạo và cung cấp cặp khóa bất đối xứng phục vụ trao đổi khóa phiên.

\textbf{Các chức năng con:}
\begin{itemize}
	\item \textbf{Tạo cặp khóa RSA 2048 bit:} Sinh public key và private key bằng bộ sinh số ngẫu nhiên an toàn.
	\item \textbf{Xuất khóa định dạng PEM:} Chuyển khóa về chuỗi PEM để dễ truyền và tải xuống.
	\item \textbf{Quản lý metadata khóa:} Chỉ lưu metadata public key (nếu cần); private key không lưu trên CSDL.
\end{itemize}

\subsection{Chức năng 2: Mã hóa tệp (Sender)}
\textbf{Mô tả:} Người gửi tải tệp và public key của người nhận, hệ thống tạo payload mã hóa.

\textbf{Các chức năng con:}
\begin{itemize}
	\item \textbf{Tạo session key:} Tạo khóa AES 32 byte và IV ngẫu nhiên.
	\item \textbf{Mã hóa tệp bằng AES-GCM:} Đầu vào là fileBuffer, đầu ra gồm ciphertext và tag.
	\item \textbf{Mã hóa session key bằng RSA-OAEP:} Dùng public key người nhận để mã hóa khóa AES.
	\item \textbf{Đóng gói payload:} Tạo tệp .hybrid gồm encrypted\_key, iv, ciphertext, tag và metadata tệp.
\end{itemize}

\subsection{Chức năng 3: Giải mã tệp (Receiver)}
\textbf{Mô tả:} Người nhận cung cấp payload .hybrid và private key để khôi phục tệp gốc.

\textbf{Các chức năng con:}
\begin{itemize}
	\item \textbf{Phân tích payload:} Tách các trường dữ liệu bắt buộc và kiểm tra định dạng.
	\item \textbf{Giải mã session key:} Dùng private key RSA giải mã encrypted\_key.
	\item \textbf{Giải mã ciphertext bằng AES-GCM:} Sử dụng session key + IV và xác minh tag.
	\item \textbf{Xuất tệp gốc:} Trả về fileBuffer và tên tệp ban đầu.
\end{itemize}

\subsection{Chức năng 4: Xử lý payload}
\textbf{Mô tả:} Chuẩn hóa dữ liệu trung gian để truyền, lưu tạm, và tái sử dụng.

\textbf{Các chức năng con:}
\begin{itemize}
	\item \textbf{Đóng gói payload JSON:} Mã hóa Base64 cho encrypted\_key, iv, ciphertext, tag.
	\item \textbf{Ràng buộc kích thước và validate:} Giới hạn dung lượng tệp, kiểm tra trường bắt buộc.
	\item \textbf{Ghi log metadata:} Lưu thông tin tệp (tên, kích thước) và trạng thái encrypt/decrypt.
\end{itemize}

\section{Dòng dữ liệu giữa các chức năng}
Luồng dữ liệu của hệ thống được thực hiện theo trình tự:
\begin{itemize}
	\item \textbf{Bước 1 -- Key Management:} Đầu vào là yêu cầu tạo khóa; đầu ra là public\_key, private\_key.
	\item \textbf{Bước 2 -- File Encryption:} Đầu vào là fileBuffer + public\_key; đầu ra là payload .hybrid.
	\item \textbf{Bước 3 -- Payload Handling:} Payload được ghi/tải/xuất dưới dạng tệp nhị phân hoặc JSON.
	\item \textbf{Bước 4 -- File Decryption:} Đầu vào là payload + private\_key; đầu ra là file gốc.
	\item \textbf{Bước 5 -- Metadata Logging:} Mỗi bước encrypt/decrypt cập nhật bảng files và logs.
\end{itemize}

\section{Sơ đồ dòng dữ liệu (DFD)}
\begin{figure}[H]
\centering
\begin{tikzpicture}[
	node distance=1.5cm and 1.8cm,
	>=Stealth,
	font=\small,
	process/.style={rectangle, draw, text width=2.4cm, minimum height=1.0cm, align=center},
	actor/.style={rectangle, draw, rounded corners, text width=1.8cm, minimum height=0.9cm, align=center},
	store/.style={cylinder, draw, shape border rotate=90, aspect=0.28, minimum height=1.0cm, minimum width=2.2cm, align=center}
]

\node[actor] (sender) {Người gửi};
\node[process, right=of sender] (aes) {Mã hóa\\AES-256-GCM};
\node[process, above=of aes] (rsaenc) {Mã hóa khóa phiên\\RSA-OAEP};
\node[process, right=of aes] (payload) {Tạo payload\\(.hybrid)};
\node[store, below=of payload] (db) {Kho lưu trữ};
\node[actor, right=of payload] (receiver) {Người nhận};
\node[process, below=of receiver] (dec) {Giải mã\\RSA + AES-GCM};

\draw[->] (sender) -- node[above] {tệp} (aes);
\draw[->] (sender) |- node[left,pos=0.25] {public key} (rsaenc);
\draw[->] (aes) -- node[above] {ciphertext, tag} (payload);
\draw[->] (rsaenc) -- node[above] {encrypted key} (payload);
\draw[->] (payload) -- node[above] {.hybrid} (receiver);
\draw[->] (receiver) -- node[right] {private key, payload} (dec);
\draw[->] (payload) -- node[right] {metadata} (db);
\draw[->] (dec) |- node[left,pos=0.28] {log} (db);

\end{tikzpicture}
\caption{Data Flow Diagram cho hệ thống truyền tệp mã hóa lai}
\end{figure}

\section{Mô hình quan hệ thực thể (ERD)}
\begin{figure}[H]
\centering
\begin{tikzpicture}[
	node distance=1.8cm and 2.2cm,
	>=Stealth,
	entity/.style={rectangle, draw, minimum width=3.2cm, minimum height=1.2cm, align=center}
]

\node[entity] (users) {Users\\(user\_id, name, email)};
\node[entity, right=of users] (keys) {Keys\\(key\_id, user\_id, public\_key)};
\node[entity, below=of users] (files) {Files\\(file\_id, filename, size, created\_at)};
\node[entity, right=of files] (transfers) {Transfers\\(transfer\_id, sender\_id, receiver\_id, file\_id, status)};

\draw[->] (users) -- node[above] {1--N} (keys);
\draw[->] (users) -- node[left] {1--N} (files);
\draw[->] (users) -- node[right,pos=0.5] {1--N (sender/receiver)} (transfers);
\draw[->] (files) -- node[above] {1--N} (transfers);
\draw[->] (keys) -- node[right] {1--1 active key} (transfers);

\end{tikzpicture}
\caption{ERD cho mô hình metadata hệ thống}
\end{figure}

\section{Mô tả bảng dữ liệu}
\begin{longtable}{|p{3cm}|p{4.2cm}|p{7cm}|}
\hline
\textbf{Tên bảng} & \textbf{Thuộc tính} & \textbf{Mô tả} \\
\hline
Users & user\_id (PK), name, email, role & Lưu thông tin người gửi/nhận trong hệ thống. \\
\hline
Keys & key\_id (PK), user\_id (FK), public\_key, created\_at & Quản lý public key của người dùng; private key không lưu trong CSDL. \\
\hline
Files & file\_id (PK), filename, size, created\_at & Lưu metadata tệp đã mã hóa/giải mã, không lưu nội dung rõ. \\
\hline
Transfers & transfer\_id (PK), sender\_id (FK), receiver\_id (FK), file\_id (FK), status, timestamp & Theo dõi phiên truyền tệp, kết quả encrypt/decrypt và trạng thái xử lý. \\
\hline
\end{longtable}

\section{Kết luận}
Hệ thống truyền tệp mã hóa lai đáp ứng yêu cầu bảo mật hiện đại thông qua việc kết hợp AES-256-GCM và RSA-OAEP. Cách tiếp cận này bảo đảm hiệu năng cho tệp lớn, đồng thời duy trì tính bí mật và toàn vẹn dữ liệu. Thiết kế phân lớp chức năng rõ ràng, kèm theo DFD và ERD, tạo nên nền tảng tốt cho việc mở rộng trong thực tế, như quản lý người dùng, kiểm toán log và tự động hóa quy trình truyền nhận tệp.

\end{document}
